%%%%%%%% GenBio @ ICML 2026 — Workshop Draft v4 %%%%%%%%%%%%%%%%%
% v4: Replace Fourier with 2-Component; add β_s literature match;
%     model-dependent correction feasibility; enriched discussion

\documentclass{article}

\usepackage{microtype}
\usepackage{graphicx}
\usepackage{subcaption}
\usepackage{booktabs}
\usepackage{hyperref}

\newcommand{\theHalgorithm}{\arabic{algorithm}}

\usepackage{icml2026}

\usepackage{amsmath}
\usepackage{amssymb}
\usepackage{mathtools}
\usepackage{xcolor}
\usepackage{multirow}

\usepackage[capitalize,noabbrev]{cleveref}

% Placeholder command for figures not yet created
\newcommand{\figplaceholder}[2]{%
  \begin{center}
    \fbox{\parbox{#1}{\centering\vspace{1.5em}\textcolor{gray}{\textbf{[FIGURE PLACEHOLDER]}\\#2}\vspace{1.5em}}}
  \end{center}
}

\icmltitlerunning{Personalized Distortion Modeling in High-Dimensional Neural Color Representations}

\begin{document}

\twocolumn[
  \icmltitle{Inferring Subject-Specific Distortion Fields from \\
    High-Dimensional Neural Color Representations}

  \icmlsetsymbol{equal}{*}

  \begin{icmlauthorlist}
    \icmlauthor{Anonymous Authors}{anon}
  \end{icmlauthorlist}

  \icmlaffiliation{anon}{Anonymous Institution}

  \icmlcorrespondingauthor{Anonymous Author}{anonymous@example.com}

  \icmlkeywords{color vision deficiency, forward encoding, fMRI, personalized modeling, representational geometry}

  \vskip 0.3in
]

\printAffiliationsAndNotice{}

%==============================================================================
\begin{abstract}
%==============================================================================
Color vision deficiency (CVD) distorts the continuous geometry of cortical color representations, but characterizing this distortion from high-dimensional fMRI signals and mapping it to low-dimensional mechanistic parameters remains challenging.
We introduce a framework for fitting parsimonious, retina-inspired distortion models (1--2 free parameters) to individual fMRI-measured hue interpolation profiles in visual cortex.
Applied to two CVD and one control participant, the framework reveals that each CVD case is best captured by a \emph{different} model---a 1-parameter retinal cone-shift for moderate protanomaly and a 2-parameter retinal--cortical model for mild deuteranomaly---while the control is correctly rejected by all models.
A complementary cortical-level model recovers an S-cone compensatory rotation of $\sim$21.5$^\circ$, matching the 21.4$^\circ$ independently reported from behavioral hue-scaling, and is the only model validated across both interpolation and pairwise-distance criteria.
The fitted models further predict whether stimulus-space correction is feasible: the degree of geometric collapse in the distortion field determines correction limits, with model choice itself shaping what is deemed correctable.
\end{abstract}

%==============================================================================
\section{Introduction}
\label{sec:intro}
%==============================================================================

Color vision deficiency (CVD) arises from altered spectral sensitivity of retinal cone photoreceptors and affects approximately 8\% of males worldwide \cite{placeholder_machado2009}.
Current assistive approaches, including commercial spectral-notch filters, apply a fixed, generic correction.
Recent evidence indicates that such filters can alter color \emph{appearance} but have minimal effect on color \emph{discrimination} at threshold \cite{placeholder_somers2024}, suggesting that effective correction may require individually-tailored strategies informed by each person's neural color geometry.

A central challenge is that CVD-related neural distortion is not uniform across individuals.
Even within a single diagnostic category (e.g., deuteranomaly), the degree of cortical compensation varies substantially \cite{placeholder_tregillus2021}, and the functional deficit that matters for continuous color tasks---the ability to interpolate between adjacent hues, not merely classify them---differs in ways that categorical diagnosis cannot capture.
Characterizing this distortion requires moving from classification-based neural decoding to \emph{geometry-sensitive} evaluation of high-dimensional fMRI representations: asking not whether the brain can tell two colors apart, but whether the continuous structure between them is preserved.

Here we propose a framework for inferring subject-specific distortion fields from fMRI data, casting CVD as a case study in recovering low-dimensional distortion structure from complex representational geometry.
We fit parsimonious mechanistic forward models to each CVD participant's neural interpolation profile, as measured by a geometry-sensitive leave-one-color-out (LOCO) cross-validation metric.
Our contributions are:
\begin{enumerate}
    \item We identify subject-specific interpolation failure as a geometry-sensitive phenotype that dissociates from classification accuracy, with the effect most robustly captured in hV4.
    \item We show that different CVD participants are best captured by different parsimonious models (1--2 DOF), with successful specificity rejection of a normal control.
    A complementary cortical-level model recovers an S-cone compensatory rotation converging with independently measured behavioral data.
    \item We demonstrate that correction feasibility depends on the degree of geometric collapse induced by the fitted model, and that model choice itself determines what is deemed correctable---retinal models predict irreversible collapse for moderate CVD, while cortical-level models preserve bijectivity.
\end{enumerate}

%==============================================================================
\section{Methods}
\label{sec:methods}
%==============================================================================

\subsection{Participants and Stimuli}

Ten participants were included in the analysis (7 healthy controls [HC], 3 females, age $22.7 \pm 2.5$ years; 3 color vision deficient [CVD], all male, including one deuteranope, one protanomalous, and one near-normal deutan serving as a specificity control).
Color vision was assessed using the 14-plate Ishihara test \cite{placeholder_ishihara1917}.
Given the small CVD sample, all CVD results are interpreted as individual case demonstrations, not population-level effects \cite{placeholder_crawford1998}.

Stimuli consisted of eight colors equally spaced from 0$^\circ$ to 315$^\circ$ in CIE~$L^*a^*b^*$ space, with lightness fixed at $L^* = 75$ and chroma distance of 40 from the achromatic point.
Participants performed a modified Rapid Serial Visual Presentation (RSVP) detection task \cite{placeholder_brouwer2009} across six runs per session, maintaining fixation while viewing 1.5~s color patches presented in optimized order.
Functional images were acquired on a Siemens 3T scanner (TR\,=\,1.5~s, 2~mm isotropic voxels, 24 oblique slices perpendicular to the calcarine sulcus), preprocessed with fMRIPrep, and projected to retinotopically-defined ROIs (V1, V2, hV4) using the Wang probabilistic atlas at 50\% threshold \cite{placeholder_wang2015}.
Response amplitudes were estimated using a two-stage FIR-then-GLM procedure \cite{placeholder_brouwer2009} and Procrustes-aligned across runs within each participant to improve cross-run consistency.

\subsection{Neural Phenotype: Interpolation Vulnerability}

We modeled hue-selective neural populations using half-wave rectified and squared sinusoidal basis functions \cite{placeholder_brouwer2009}: each of $K$ hypothetical channels was tuned to a preferred hue $\theta_k$, with response $r_k(\theta) = \max(0, \cos(\theta - \theta_k))^2$.
The weight matrix $W$ mapping channel outputs to voxel responses was estimated via ridge regression with generalized cross-validation (GCV) for regularization \cite{placeholder_golub1979}.

The target phenotype is the \emph{LOCO vulnerability profile}, a geometry-sensitive metric of representational structure.
Leave-one-color-out (LOCO) cross-validation tests whether the model can interpolate to novel hues not seen during training, assessing the preservation of continuous color-space geometry in the high-dimensional voxel pattern space: $W$ is trained on 7 of 8 colors and voxel responses to the held-out color are predicted.
The per-color prediction error defines an 8-element vulnerability vector $\mathbf{v} \in \mathbb{R}^8$ for each participant, capturing which hues are poorly interpolated from their neighbors---a structural deficit invisible to classification-based decoders.

This phenotype is specific to continuous geometry and distinct from color discrimination.
Leave-one-\emph{run}-out (LORO) cross-validation, which tests whether individual colors can be reliably distinguished across repetitions, shows no HC--CVD difference ($p = 0.668$), while LOCO interpolation error is significantly elevated in CVD.
Among the ROIs tested, the effect was most robust in hV4 ($p = 0.017$; $K{=}3$), consistent with hV4's established role in perceptual color representation \cite{placeholder_brouwer2009, placeholder_bannert2018}.
We therefore focus distortion fitting on hV4, with V1 analyses reported as cross-ROI validation in \cref{app:supporting}.

\subsection{Distortion Models}

We fit three forward models of hue distortion (\cref{fig:main}a), each parameterizing how a physical stimulus angle $\theta$ is transformed into a distorted angle $\theta'$:

\paragraph{Machado cone shift (1 DOF).}
A physiological model \cite{placeholder_machado2009} in which the anomalous cone's spectral peak is shifted by $\Delta\lambda \in [0, 20]$~nm.
This retinal-level model predicts hue distortion through altered L--M cone opponency.

\paragraph{Retinal + cortical gain (R+C; 2 DOF).}
Extends the cone shift with a cortical opponent-channel gain $g$:
\begin{equation}
rg' = rg_{\text{base}} + (1{+}g) \cdot (rg_{\text{ret}} - rg_{\text{base}})
\label{eq:rc}
\end{equation}
where $g{=}0$ reduces to pure Machado, $g{=}{-}1$ is exact cortical compensation, and $g{<}{-}1$ is overcompensation \cite{placeholder_tregillus2021}.

\paragraph{2-component angular dilation (2 DOF).}
An effective cortical-level model that parameterizes hue distortion as angular dilation along two opponent axes:
\begin{equation}
\theta'(c) = \theta(c) + \beta_s \cos(\theta(c) {-} 90^\circ) + \beta_c \cos(\theta(c) {-} \theta_{\text{conf}})
\label{eq:2comp}
\end{equation}
where $\beta_s$ captures S-cone axis expansion (compensatory upregulation of the intact S-cone pathway) and $\beta_c$ captures confusion-axis modulation specific to CVD family ($\theta_{\text{conf}} = 16^\circ$ for protan, $150^\circ$ for deutan).
Unlike the cone-level models, this operates directly in opponent hue space without modeling retinal spectral shifts.

\subsection{Fitting and Specificity}

For each model and candidate parameter(s), we simulate the LOCO vulnerability profile that a healthy observer would exhibit if viewing stimuli distorted by the model's forward map.
The fit is the Spearman rank correlation $\rho$ between simulated and observed vulnerability vectors, optimized over a grid ($\Delta\lambda$: 41 points; $g$: 25 points; $\beta_s, \beta_c$: $26 \times 51 = 1{,}326$ points).

Statistical significance is assessed by exact permutation of the 8-color labels ($8! = 40{,}320$ permutations).
A valid model must satisfy two criteria:
\begin{enumerate}
    \item \textbf{Sensitivity}: significant fit ($p < 0.05$) for CVD participants.
    \item \textbf{Specificity}: non-significant fit for the normal control (sub-10).
\end{enumerate}

Additionally, we evaluate a pairwise-distance criterion ($\Delta$RDM) independently on V1, fitting the 2-component model to the difference between CVD and mean-HC representational distance matrices in voxel space (crossnobis cosine similarity; 8! exact permutations + bootstrap CI).
A model validated under both LOCO (per-color interpolation accuracy) and $\Delta$RDM (pairwise geometric structure) provides stronger evidence than either criterion alone.

%==============================================================================
\section{Results}
\label{sec:results}
%==============================================================================

\subsection{Subject-Specific Model Selection}

\cref{tab:main} and \cref{fig:main} present the central result: distinct CVD participants require distinct distortion models, while the normal control yields no significant fit.

\begin{table}[t]
  \caption{Subject-specific distortion model fits (hV4 LOCO). Each CVD subject is best explained by a different parsimonious model. The normal control is correctly rejected by all models.}
  \label{tab:main}
  \begin{center}
    \begin{small}
      \begin{tabular}{llcccc}
        \toprule
        Subject & Model & DOF & $\rho$ & $p_{\text{perm}}$ \\
        \midrule
        Sub-08 (deutan) & R+C & 2 & 0.857 & \textbf{0.005**} \\
         & 2-Comp & 2 & 0.881 & \textbf{0.004**} \\
        Sub-09 (protan) & Machado & 1 & 0.762 & \textbf{0.018*} \\
         & 2-Comp & 2 & 0.690 & \textbf{0.035*} \\
        Sub-10 (control) & --- & --- & $-0.048$ & 0.559 \\
        \bottomrule
      \end{tabular}
    \end{small}
  \end{center}
  \vskip -0.15in
\end{table}

\textbf{Sub-09 (moderate protanomaly).}
The 1-DOF Machado model with $\Delta\lambda = 13.5$~nm---within the established range for moderate--severe protanomaly \cite{placeholder_machado2009}---captures the interpolation profile at $\rho = 0.762$, $p = 0.018$.
Adding a cortical gain parameter does not improve fit ($g = 0.0$ at optimum), indicating that a purely retinal model is sufficient.
The 2-component model also reaches significance ($p = 0.035$) but with lower $\rho$, confirming that for this participant the retinal-level account is most parsimonious.

\textbf{Sub-08 (mild deuteranomaly).}
The 1-DOF Machado model is only trend-level ($p = 0.058$).
The 2-DOF R+C model, with a small retinal shift ($\Delta\lambda = 2.0$~nm) and a cortical gain term ($g = +2.25$), achieves $\rho = 0.857$, $p = 0.005$.
The 2-component model achieves the highest fit of any model tested ($\rho = 0.881$, $p = 0.004$) with $\beta_s = 38^\circ$ and $\beta_c = -14^\circ$.
Both 2-DOF models substantially outperform the retinal-only account, confirming that the additional parameter is essential for this participant.

\textbf{Sub-10 (normal control).}
All models yield non-significant fits ($p \geq 0.559$ for cone-level; $p = 0.058$ marginal for 2-component), confirming specificity.

\begin{figure}[t]
  \vskip 0.2in
  \figplaceholder{\columnwidth}{%
    \textbf{Figure 1: Subject-Specific Distortion Profiles}\\[0.3em]
    \small \textbf{(a)} Model hierarchy: Machado (1 DOF, retinal) $\to$ R+C (2 DOF, retinal+cortical) $\to$ 2-Component (2 DOF, cortical-effective). Each operates at a different level of the visual processing stream.\\[0.3em]
    \textbf{(b--d)} 8-color LOCO vulnerability profiles.
    Gray bars: observed CVD profile.
    Colored lines: best-model fits.
    \textbf{(b)} Sub-08: R+C and 2-Comp fits ($p{=}0.005$, $p{=}0.004$). \textbf{(c)} Sub-09: Machado fit ($\rho{=}0.762$, $p{=}0.018$). \textbf{(d)} Sub-10: flat profile, $p{=}0.559$ (specificity pass).
  }
  \caption{The central result. Each CVD participant exhibits a distinct vulnerability profile (gray bars) in the high-dimensional voxel representation, captured by a different parsimonious model (colored lines). The normal control shows no structured vulnerability. Panel (a) illustrates the model hierarchy from retinal to cortical levels.}
  \label{fig:main}
  \vskip -0.1in
\end{figure}

\subsection{Fitted Parameters and Cross-Criterion Validation}

The fitted retinal parameters fall within established physiological ranges: sub-08's $\Delta\lambda = 2.0$~nm matches very mild deuteranomaly, and sub-09's $\Delta\lambda = 13.5$~nm matches moderate protanomaly \cite{placeholder_machado2009}.

Sub-08's R+C cortical gain ($g = +2.25$) is not straightforward to interpret.
One possibility is that the mild retinal deficit produces a small spectral signal that cortical processing amplifies into a measurable interpolation deficit---an effect invisible under retinal-only modeling.
However, $g$ could alternatively absorb unmodeled variance.
We treat this as a provisional effective description, noting that the specific sign (amplification rather than compensation) lacks direct precedent \cite{placeholder_tregillus2021}.

The 2-component model provides an independent cross-criterion check.
Under the $\Delta$RDM criterion in V1 (pairwise distance structure, not per-color accuracy), bootstrap confidence intervals for the S-cone expansion parameter are $\beta_s = 20.0^\circ \pm 8.0^\circ$ (sub-08) and $\beta_s = 23.0^\circ \pm 10.2^\circ$ (sub-09), both excluding zero.
The cross-subject mean $\beta_s \approx 21.5^\circ$ matches the $21.4^\circ$ B-Y axis rotation independently measured via behavioral hue-scaling \cite{placeholder_emery2021}---convergent evidence from an entirely different method, criterion (geometric distance vs.\ perceptual scaling), and brain region (V1 vs.\ behavior).
This convergence suggests that the S-cone compensatory component reflects a genuine biological signal rather than an artifact of the fitting procedure.

The confusion-axis parameter $\beta_c$ shows family-appropriate behavior: significantly negative for the deutan ($-18^\circ \pm 6^\circ$, CI excluding 0) and non-significant for the protan ($+3^\circ \pm 2^\circ$, CI including 0), consistent with the larger retinal shift in protanomaly already capturing confusion-axis structure.

The 2-component model is the only model validated under both the LOCO criterion (sub-08 hV4 $p = 0.004$; sub-09 hV4 $p = 0.035$) and the $\Delta$RDM criterion (sub-09 V1 $p = 0.007$; sub-08 V1 bootstrap CI excluding 0).
No other model achieves significance under both criteria for both CVD participants (\cref{tab:full}).

\subsection{Model-Dependent Correction Limits}

Given a fitted forward model $D\!: \theta \to \theta'$, we ask whether an inverse correction exists---i.e., whether modified stimulus angles $\theta^*$ can be found such that $D(\theta^*) = \theta_{\text{target}}$ for all 8 colors.

\textbf{Sub-08 (mild):} Both the R+C and 2-component forward maps are bijective, and numerical pre-image search yields exact solutions ($< 0.001^\circ$ residual) for all 8 colors under both models---producing 8-point stimulus correction lookup tables (\cref{fig:feasibility}a).
The two models prescribe different correction directions (cosine between correction vectors = $-0.18$; sign agreement on 3/8 colors), reflecting their distinct mechanistic assumptions.
Which correction actually improves perception requires behavioral validation.

\textbf{Sub-09 (moderate):} The answer depends on model choice.
Under the best-fitting \emph{retinal} model (Machado, $\Delta\lambda = 13.5$~nm), the forward map compresses the 360$^\circ$ opponent arc to approximately 96$^\circ$, merging green, cyan, and blue into a single perceived angle ($\sim$282$^\circ$).
Exact inversion is structurally impossible for 4/8 colors (\cref{fig:feasibility}b).
A discriminability-oriented fallback achieves 5.6$\times$ improvement but reaches only 48\% of healthy-observer separation.

However, under the \emph{cortical-level} 2-component model ($\beta_s = 6^\circ$, $\beta_c = -22^\circ$, $p = 0.035$), the forward map is an angular dilation without arc compression and therefore remains bijective.
All 8 colors admit exact pre-image solutions (mean $|\delta| = 20.1^\circ$, max $= 48.1^\circ$).
This contrast reveals that correction feasibility is not solely a property of severity: it depends on which level of the visual hierarchy the distortion model targets.
Retinal models can exhibit irreversible spectral collapse; cortical-level models, operating in opponent hue space, preserve bijectivity by construction.

\begin{figure}[t]
  \vskip 0.2in
  \figplaceholder{\columnwidth}{%
    \textbf{Figure 2: Model-Dependent Correction Limits}\\[0.3em]
    \small \textbf{(a)} Sub-08 (mild): both R+C and 2-Component admit exact pre-image for 8/8 colors ($<0.001^\circ$ error), but correction directions differ (cosine = $-0.18$).\\
    \textbf{(b)} Sub-09 (moderate): Machado compresses 360$^\circ$ $\to$ $\sim$96$^\circ$ (exact inversion impossible, 4/8 only). 2-Component preserves bijectivity $\to$ 8/8 exact. Correction feasibility depends on model choice.
  }
  \caption{Correction feasibility depends on both severity and model level. Retinal models can exhibit irreversible arc compression; cortical-level angular dilation preserves bijectivity. For the moderate CVD participant, model choice determines whether exact correction is structurally possible.}
  \label{fig:feasibility}
  \vskip -0.1in
\end{figure}

%==============================================================================
\section{Discussion}
\label{sec:discussion}
%==============================================================================

\paragraph{Inferring distortion from complex signals.}
The core methodological contribution is that subject-specific, low-dimensional distortion fields (1--2 parameters) can be inferred from high-dimensional fMRI voxel patterns.
The LOCO vulnerability profile serves as the bridge: it compresses the high-dimensional representation into a geometry-sensitive 8-element summary that captures continuous structure rather than categorical boundaries.
This approach may generalize beyond CVD to other conditions where sensory representations are distorted rather than absent.

\paragraph{S-cone compensation as a convergent biological signal.}
The 2-component model's S-cone expansion parameter ($\beta_s \approx 21.5^\circ$) converges with Emery et al.'s \cite{placeholder_emery2021} independently measured 21.4$^\circ$ B-Y rotation from behavioral hue-scaling---a striking match across methods (fMRI $\Delta$RDM vs.\ psychophysics), brain regions (V1 vs.\ behavior), and CVD types (deutan and protan).
This convergence is consistent with a compensatory upregulation of the intact S-cone pathway that is shared across CVD subtypes, as has been hypothesized from adaptation studies \cite{placeholder_tregillus2021, placeholder_boehm2014}.
That the same parameter emerges from both per-color interpolation fitting (LOCO) and pairwise distance fitting ($\Delta$RDM) further supports its robustness, though the specific $\beta_s$ values differ across criteria (LOCO: 38$^\circ$ for sub-08; $\Delta$RDM: 20$^\circ$ $\pm$ 8$^\circ$), likely reflecting different aspects of the same underlying distortion.

\paragraph{Model choice shapes intervention design.}
An unexpected finding is that correction feasibility is not determined by CVD severity alone but depends on which level of the visual hierarchy the forward model targets.
A retinal-level model predicts irreversible spectral collapse for moderate protanomaly; a cortical-level model applied to the same participant preserves bijectivity.
This is not merely an academic distinction: it determines whether a stimulus-space correction lookup table can be constructed.
Resolving which model is correct for a given individual is thus practically consequential, and requires behavioral validation of the model-prescribed corrections.

\paragraph{Why different models for different subjects?}
Sub-09's profile is captured by pure cone shift ($g = 0$) while sub-08 requires the additional cortical gain term.
This heterogeneity suggests that the neural consequences of CVD cannot be reduced to diagnostic category alone, but instead reflect person-specific variation in retinal severity and downstream processing that is only visible under individualized modeling.

\paragraph{Limitations.}
Three CVD participants are insufficient for population-level claims; this work is a proof-of-concept.
The fMRI requirement limits near-term scalability.
The 2-component model's marginal specificity (sub-10 $p = 0.058$) raises the possibility of overfitting with 1,326 grid points on 8 colors, though the $\Delta$RDM validation in V1 provides independent confirmation.
The interpolation phenotype depends on forward encoding model specification (basis type, number of channels), though prior analyses confirm robustness across reasonable parameter choices.

\paragraph{Relation to prior work.}
Our approach extends Brouwer \& Heeger's \cite{placeholder_brouwer2009} forward encoding framework---which demonstrated LOCO reconstruction in hV4 for healthy observers---to CVD with individual-level distortion fitting.
The distinction between preserved classification and impaired interpolation is consistent with reports that CVD individuals maintain categorical but not continuous color representations \cite{placeholder_bannert2018}.
The cortical compensation parameters relate to accumulating evidence for long-term neural adaptation in anomalous trichromacy \cite{placeholder_tregillus2021, placeholder_emery2021}.

%==============================================================================
\section{Conclusion}
\label{sec:conclusion}
%==============================================================================

We present a proof-of-concept framework for inferring subject-specific distortion fields from high-dimensional neural representations in color vision deficiency.
The framework identifies individual distortion profiles with parsimonious models, recovers a compensatory S-cone rotation that converges with independent behavioral measurements, and reveals that correction feasibility depends on both severity and the level of the visual hierarchy targeted by the forward model.
The critical next step is behavioral validation: determining whether the model-prescribed corrections actually improve color discrimination in the fitted individuals.

%==============================================================================
\section*{Impact Statement}
%==============================================================================

This work develops personalized neural models of color vision deficiency with the goal of informing assistive technology design.
Potential positive impacts include individually-tailored color corrections for digital displays and accessibility tools.
The framework currently requires fMRI data, limiting immediate clinical deployment.
Behavioral validation and informed consent are essential prerequisites before any correction is applied, as a misspecified model could impair rather than improve visual function.

%==============================================================================
\bibliography{genbio_draft}
\bibliographystyle{icml2026}

%==============================================================================
%==============================================================================
\newpage
\appendix
\onecolumn

\section{Supporting Analyses}
\label{app:supporting}

\subsection{Cross-ROI and Cross-Criterion Validation}

Beyond the primary hV4 LOCO fitting (main text), the 2-component model was independently validated in V1 using pairwise representational distance matrices ($\Delta$RDM).
Sub-09 achieved crossnobis cosine $p = 0.007$ with $\beta_s = 23.0^\circ \pm 10.2^\circ$ (bootstrap CI excluding 0).
Sub-08 V1 LOCO also reached $p = 0.001$ ($\beta_s = 50^\circ$, $\beta_c = -14^\circ$)---the strongest LOCO result in the pipeline---confirming that the hV4-based finding generalizes across ROIs.

The cross-subject mean $\beta_s \approx 21.5^\circ$ from V1 $\Delta$RDM matches the $21.4^\circ$ B-Y axis rotation independently reported from behavioral hue-scaling \cite{placeholder_emery2021}, providing convergent evidence from entirely different fitting criteria and brain regions.

\subsection{Model Comparison: Full Table}

\begin{table}[h]
  \caption{Full model comparison across hV4 LOCO, V1 LOCO, and V1 $\Delta$RDM criteria. The 2-component model is the only model achieving significance under both LOCO and $\Delta$RDM for both CVD participants.}
  \label{tab:full}
  \begin{center}
    \begin{small}
      \begin{tabular}{llcccccc}
        \toprule
        Subject & Model & DOF & hV4 LOCO $\rho$ & hV4 LOCO $p$ & V1 LOCO $p$ & V1 $\Delta$RDM & Notes \\
        \midrule
        \multirow{3}{*}{Sub-08} & Machado & 1 & 0.619 & 0.058 & 0.033* & anti-corr. & Trending (hV4) \\
        & R+C & 2 & 0.857 & \textbf{0.005**} & 0.047* & 0.179 & hV4 primary \\
        & 2-Comp & 2 & \textbf{0.881} & \textbf{0.004**} & \textbf{0.001***} & CI excl.\ 0 & Dual-validated \\
        \midrule
        \multirow{3}{*}{Sub-09} & Machado & 1 & \textbf{0.762} & \textbf{0.018*} & 0.112 & $+0.091$ & hV4 primary \\
        & R+C & 2 & 0.762 & 0.018 & --- & 0.026* & $g{=}0$ (= Machado) \\
        & 2-Comp & 2 & 0.690 & 0.035* & 0.018* & \textbf{0.007***} & Dual-validated \\
        \midrule
        Sub-10 & All & --- & $-0.048$ & 0.559 & 0.058$^\dagger$ & --- & Specificity \\
        \bottomrule
        \multicolumn{8}{l}{\small $^\dagger$ 2-component only; cone-level models all NS. Marginal---see Limitations.} \\
      \end{tabular}
    \end{small}
  \end{center}
\end{table}

\subsection{Correction Feasibility Details}

\paragraph{Sub-08 (R+C pre-image).}
All 8 colors admit exact pre-image solutions (mean $|\delta| = 23.5^\circ$, max $= 42.9^\circ$, residual $< 0.001^\circ$).

\paragraph{Sub-08 (2-component pre-image).}
All 8 colors admit exact pre-image solutions (mean $|\delta| = 46.3^\circ$, max $= 104.2^\circ$, residual $< 0.001^\circ$).
The two models prescribe substantially different correction magnitudes (R+C mean 23.5$^\circ$ vs.\ 2-component mean 46.3$^\circ$) and partially opposite directions (cosine = $-0.18$, sign agreement 3/8).

\paragraph{Sub-09 (Machado pre-image).}
At $\Delta\lambda = 13.5$~nm, the opponent-space arc compresses from 360$^\circ$ to $\sim$96$^\circ$.
Colors c4 (green), c5 (cyan), and c6 (blue) converge to a single perceived angle ($\sim$282$^\circ$).
Exact pre-image: 4/8 colors.
Separation-optimized fallback: min separation $1.03^\circ \to 5.76^\circ$ ($5.6\times$), mean separation $39.3^\circ \to 24.3^\circ$ (healthy: 70.7$^\circ$).

\paragraph{Sub-09 (2-component pre-image).}
All 8 colors admit exact pre-image solutions (mean $|\delta| = 20.1^\circ$, max $= 48.1^\circ$, residual $< 0.001^\circ$).
The cortical-level angular dilation, unlike retinal spectral shift, does not compress the opponent arc and therefore preserves bijectivity.

\end{document}
